\documentclass[a4paper,12pt]{article}
\usepackage[utf8]{inputenc}
\usepackage[T1]{fontenc}
\usepackage[ngerman]{babel}
\usepackage{amsmath}
\usepackage{amssymb}
\usepackage{enumitem}
\usepackage{geometry}
\geometry{a4paper, margin=2.5cm}
\usepackage{parskip}

\title{Einsendeaufgaben -- Aufgabe 4.5\\[0.5em]
\large Modul 61112: Lineare Algebra}
\author{}
\date{}

\begin{document}
\maketitle

\section*{Aufgabe 4.5}

\begin{enumerate}[label=\alph*)]
  \item
  $T$ ist eine lineare Abbildung, demnach ein Endomorphismus, denn:

  \medskip

  $\forall A \in M_{33}(K) \, \forall a \in K: \ T(aA) = a\,T(A)$
  \[
    T(aA) = (aA)^T = a A^T = a\,T(A)
  \]

  $\forall A, B \in M_{33}(K): \ T(A+B) = T(A) + T(B)$
  \[
    T(A+B) = (A+B)^T = A^T + B^T = T(A) + T(B)
  \]

  \item
  Wir wählen die geordnete Basis $B$ bestehend aus den $3 \times 3$-Elementarmatrizen
  aus:
  \[
    B := (E_{11}, E_{12}, \ldots, E_{32}, E_{33})
  \]

  Es gilt $T(E_{ij}) = E_{ji}$ für alle $i, j \in \{1, 2, 3\}$. Wir verwenden die
  Leibniz-Formel:
  \begin{align*}
    \det(T) &= \det({}_B M_B(T)) \\
            &= \sum_{\sigma \in S_9} \operatorname{sgn}(\sigma) \prod_{i=1}^{n} a_{\sigma(i), i}
  \end{align*}

  Wegen $T(E_{ij}) = E_{ji}$ ist in einer Spalte von ${}_B M_B(T)$ nur ein Eintrag $1$,
  die anderen $0$. Also gibt es nur eine Permutation $\sigma$ mit
  \[
    \prod_{i=1}^{n} a_{\sigma(i), i} \neq 0.
  \]
  Wir betrachten diese Permutation und bestimmen die Fehlstellung.

  \[
    \sigma = \begin{pmatrix}
      \overset{E_{11}}{1} & \overset{E_{12}}{2} & \overset{E_{13}}{3} & \overset{E_{21}}{4} & \overset{E_{22}}{5} & \overset{E_{23}}{6} & \overset{E_{31}}{7} & \overset{E_{32}}{8} & \overset{E_{33}}{9} \\
      1 & 4 & 7 & 2 & 5 & 8 & 3 & 6 & 9
    \end{pmatrix}
  \]

  \[
    FS(\sigma) = \{(2,4), (2,7), (3,4), (3,5), (3,7), (3,8), (5,7), (6,7), (6,8)\}
  \]

  Die Signatur ist also $\operatorname{sgn}(\sigma) = -1$. Daraus folgt
  \[
    \det(T) = \operatorname{sgn}(\sigma) = -1.
  \]
\end{enumerate}

\end{document}
